% SOURCE docs/release/2026-10-29/source/docs/writing-repair/2026-10-29/src/1de93e537048440f.tex CHAR 0
\chapter{A narrower policy interval can be more confidently wrong}
\label{gap29:policy-uncertainty}
The servicing example leaves a natural question. If a fitted transition
model uses the observed histories more efficiently than weighting, can it
give an informative interval for the gain from intervention? The answer
depends on why outcomes are missing and whether the state equation is
adequate. Resampling quantifies variation in an estimated model; it does
not establish the model's economic assumptions.

\section{Follow the customer through the calculation}
The follow-up study uses the same three-decision economy as
Chapter~\ref{gap29:servicing}. It selects a threshold policy on 2,000
training customers and evaluates that fixed choice on 1,000 independent
customers. Each of 199 bootstrap draws resamples whole evaluation
customers, including their observation histories, and refits both
transition coefficients and innovation variance.

For a plug-in value estimate $\widehat V$ and bootstrap quantiles
$q^*_{.025},q^*_{.975}$, the basic interval is
$[2\widehat V-q^*_{.975},\,2\widehat V-q^*_{.025}]$.
The improvement over never acting is calculated within the same simulated
paths and resampled customers before forming its interval. Pairing retains
the covariance between policy values. No separate censoring model is fitted:
under the baseline's conditional independence, the observed transition
regression can estimate the common transition law. The adverse dropout case
violates precisely that premise.

The selected policy is held fixed during evaluation resampling. Its target
is the value of the policy selected on that separate training sample.
Repeating selection across independent experimental replications changes
the target accordingly, and each result is compared with the reference for
its selected policy. This is different from estimating the performance of
an entire future retraining procedure.

The study uses 100 independent replications in each of three cases, with a
separate million-trajectory reference. The initial-state distribution and
reward equation are supplied by the construction. Their estimation
uncertainty is consequently absent. Numerical integration uses 8,192
scrambled low-discrepancy normal paths shared across fits and bootstrap
draws. Eight separate scrambles at the true parameters check sensitivity
to numerical integration; parameter uncertainty remains a separate source
of variation. Across those eight scrambles, standard deviations are below
0.010 reward units for values and 0.009 for improvements. This observed
sensitivity is small relative to the statistical widths reported here; it
is not a uniform integration-error bound for every fitted parameter.

\section{Precision under a maintained assumption}
For the separately selected policy, the baseline interval for improvement
has empirical coverage 94\%, with Monte Carlo standard error 2.4 percentage
points. Its average width is 0.711 reward units. The matched weighting
interval has coverage 98\% and width 6.905. The difference in width is
descriptive evidence of what the parametric restrictions buy; the sample
of replications does not establish exact calibration or general superiority.
Three bootstrap coverage decisions lie near the separate reference's
numerical uncertainty band.

A positive lower bound for improvement occurs in 79\% of baseline fitted
analyses and 2\% of the weighted analyses. This is a decision frequency,
not by itself a measure of safety. The selected list includes an unprofitable
threshold in seven of the 100 training samples. Evaluation must preserve
the selected target and the possibility of a poor selection.

\section{Why the apparently useful interval can fail}
\begin{figure}[htbp]
\centering
\includegraphics[width=.95\textwidth]{scientific-gap-closure/2026-10-29/policy-coverage.pdf}
\caption{Coverage of the separately selected policy's improvement, 100
independent replications per case. Bars are Wilson intervals for Monte Carlo
coverage uncertainty. Numerical-reference uncertainty is reported separately.
Narrow fitted intervals rely on correct observation and transition assumptions.}
\end{figure}
When dropout depends on the unseen next shock, the fitted value interval
covers its reference in none of the 100 replications. The fitted interval
for improvement covers in only 6\%, with Monte Carlo standard error
2.4 percentage points. Its bias is $+1.175$ (Monte Carlo standard error
$0.052$), and its average width is 1.140. Every fitted analysis has a
positive lower bound for improvement, including selections of an
unprofitable threshold. A compelling-looking decision can therefore be
an artefact of the observation process.

The weighted improvement interval covers in 96\% in that case, but has
bias about $+0.498$ and average width 6.348. Its value interval covers only
73\%. Cancellation and imprecision in a difference can mask biased value
estimates. The apparently satisfactory improvement coverage does not
validate the false observation model.

Omitting $X$ from the transition regression produces a different failure:
coverage for improvement falls to 61\%, with Monte Carlo standard error
4.9 percentage points. The estimated bias is $-0.253$ (Monte Carlo standard
error $0.019$), despite an average interval width of 0.689.
Resampling the fitted omission repeats it.

The calculation supports a conditional use of model-based precision and a
clear reason to test the observation and state assumptions before relying
on it. An outcome audit, better follow-up or a richer measured history can
change those assumptions' credibility. More bootstrap draws alone cannot.
These results remain a generated research comparison with known reward
and initial laws; a bank application needs its own economic specification,
sampling argument and measurement evidence.

\section{Buying the information that missing customers take away}
\label{closure30:outcome-audit}

A bank cannot recover the experience of customers who disappear merely by
fitting a more elaborate model to those who remain. It can sometimes change
the evidence. A sample of incomplete customer histories may be recovered from
archived servicing records, linked payments or carefully conducted follow-up.
The economic question is how much reliable information that extra work buys.

The preceding analysis shows why resampling a fitted model does not correct
outcome-dependent dropout. This section changes the observation design while
retaining its three-decision servicing example.
Random recovery repairs the missing-information problem under a clear
condition: the histories actually recovered must form the specified random
sample. Random invitations are insufficient if successful recovery remains
selective. Even a correct estimate can be too imprecise to support a decision.
Here a separately fitted outcome model reduces interval width substantially,
but the additional pilot information is costly and a positive policy
improvement often remains statistically unresolved.

\subsection{A known design, applied to a sequential decision}
Following up a subsample with missing outcomes is usually called double
sampling. Levis and colleagues show how noninformative successful second-stage
recovery identifies the full-data distribution despite informative original
missingness. Their Proposition~1 treats a missing outcome; Appendix~A extends
the argument to general coarsening of a record. Their mean-effect estimators
also use regression to improve efficiency under explicit conditions
\citep{levis2023double}.

This changes the information without exempting the analysis from selection
assumptions. Levis and colleagues distinguish invitations from obtaining the
missing information, and discuss nonresponse and recall error. Sun and
colleagues develop a related construction for weighted quantile treatment
effects, with conditions appropriate to quantiles \citep{sun2025double}.
A mean-reward comparison does not thereby qualify a tail-loss estimate.
Averaging conditional quantile effects does not in general recover a
population quantile effect.

The present contribution is a worked composition: recover missing pieces of
a sequential policy-value score, derive the precision consequence, and test
it in the servicing law. The two-phase idea and sequential regression
assistance are inherited methods. This does not establish an efficient
estimator for every longitudinal missing-data model.

\subsection{What must be recovered}
There are three decisions, $t=0,1,2$. Each generated customer has a fixed
characteristic $X\sim N(0,1)$ and initial state
$L_0=0.4X+\varepsilon_0$. With independent standard-normal innovations,
\begin{align*}
L_{t+1}&=0.6L_t+0.25X-0.5A_t+\varepsilon_{t+1},\\
R_t&=10-2L_{t+1}^{2}-A_t+0.5A_tX.
\end{align*}
The historical probability of action is
\[
p_t=\operatorname{clip}\{\operatorname{expit}(0.3X-0.5L_t+0.2A_{t-1}),
                       0.2,0.8\},\qquad A_{-1}=0.
\]
These are known probabilities in a generated example. Clipping belongs to
the assignment law, rather than modifying an estimated bank propensity
score. A fixed policy acting when $L_t>0.5$ is compared with never acting.
Neither is selected using evaluation outcomes.

For history $H_t$ available before action $t$, define
\[
W_t(\pi)=\prod_{s=0}^{t}
 \frac{\mathbf1\{A_s=\pi(H_s)\}}{p_s(A_s\mid H_s)},
\qquad F(\pi)=\sum_{t=0}^{2}W_t(\pi)R_t.
\]
Sequential exchangeability, consistency and positive action probabilities
give $E\{F(\pi)\}=V(\pi)$, where
$V(\pi)=E_\pi(\sum_t R_t)$. The simulation satisfies these conditions by
construction. Recovering missing records cannot make them true in a bank.

Let $C_t$ indicate that reward $t$ and its required history were routinely
observed. After dropout, later $C_t$ are zero. Split the full score:
\[
K(\pi)=\sum_t C_tW_t(\pi)R_t,\qquad
M(\pi)=\sum_t(1-C_t)W_t(\pi)R_t,\qquad F=K+M.
\]
The missing contribution may depend on unobserved outcomes. Among incomplete
records, let $B$ indicate completed recovery of the missing history, with
known probability $q>0$, independently of that history and its original
observation pattern. Complete records need no recovery and have $M=0$.
Average the following score over all initially eligible customers:
\begin{equation}
Z(\pi)=K(\pi)+\frac{B}{q}M(\pi).
\label{closure30:audit-score}
\end{equation}
Each routinely observed contribution has weight one; each recovered missing
contribution represents $1/q$ similar incomplete histories. There is no
division by the questionable routine observation probability.

Condition on the full history and routine observation pattern.
$E(B\mid\text{history},C)=q$ for incomplete records, so
$E(Z\mid\text{history},C)=F$ and $E(Z)=V(\pi)$. This requires recovery of
the quantities entering $M$: actions, states, rewards and the information
needed for assignment probabilities. A terminal balance alone may not
supply them. Recovery must measure historical outcomes faithfully and must
not change them. A follow-up intervention that changes repayment would
answer a different question.

\subsection{Identification does not imply useful precision}
Under constant independent Bernoulli recovery,
\begin{equation}
\operatorname{Var}(Z)=\operatorname{Var}(F)
                   +(q^{-1}-1)E(M^2).
\label{closure30:audit-variance}
\end{equation}
Indeed $Z-F=(B/q-1)M$ has conditional mean zero and conditional variance
$(q^{-1}-1)M^2$, so its covariance with $F$ is zero. For policy improvement,
use paired differences $F_\Delta=F(\pi)-F(0)$ and
$M_\Delta=M(\pi)-M(0)$. Evaluating both policies on the same customers
retains this pairing.

Low recovery probabilities magnify rare contributions. Incomplete customers
who match several consecutive policy actions can receive large weights.
More recovery reduces the extra variance; recovering every history still
leaves the variance of the full-history policy score.

The first routine observation mechanism continues with probability
$g_t=0.65+0.30\operatorname{expit}(0.3L_t-0.2A_t)$, independently of the
next innovation conditional on history and action. The adverse mechanism
uses $\operatorname{expit}\{\operatorname{logit}(g_t)+1.5\varepsilon_{t+1}\}$.
Weighting by the original $g_t$ in that case assumes away selection by the
unseen shock.

Each mechanism has 500 independent replications with 2,000 customers.
Recovery rates are 5\%, 10\%, 25\%, 50\% and 100\% of incomplete records,
with shared recovery draws across rates within a replication. Estimators
receive full histories only where routine observation or designated recovery
permits. The reference policy improvement is 0.75258 reward units, with
numerical Monte Carlo standard error 0.00432, from the retained
million-trajectory intervention calculation. Reward units are illustrative,
not dollars.

\begin{table}[htbp]
\centering\small
\begin{tabular}{lrrrr}
\toprule
Adverse dropout & Bias & Coverage & Width & Follow-ups\\
\midrule
Assumed routine $g_t$ & 0.755 & 0.908 & 4.506 & 0\\
Random recovery, 5\% & -0.129 & 0.944 & 11.809 & 62\\
Random recovery, 25\% & -0.052 & 0.948 & 6.047 & 308\\
Random recovery, 50\% & 0.016 & 0.930 & 4.772 & 618\\
Random recovery, 100\% & 0.029 & 0.942 & 3.933 & 1,238\\
\bottomrule
\end{tabular}
\caption{Fixed-policy improvement under outcome-dependent dropout;
500 replications, 2,000 evaluation customers. Width is the mean width of a
normal 95\% interval. Follow-ups are completed missing-history recoveries.
No pilot is used. Coverage is against the numerical reference, not a
finite-sample guarantee.}
\end{table}

The estimated bias at 25\% recovery is small relative to its Monte Carlo
standard error, 0.070. Yet mean interval width is about eight times the
improvement. Even full recovery yields a positive lower bound in only
11.6\% of replications. Under conditionally independent routine dropout,
25\% recovery gives width 5.679 and coverage 93.0\%. These results do not
establish uniform attainment of nominal coverage.

A third study makes actual recovery depend on the unseen first-missed
innovation, using
$\operatorname{expit}\{\operatorname{logit}(0.25)+0.8\varepsilon_{\rm miss}\}$,
but analyses it as constant 25\% recovery. Across 500 replications, upward
bias is 0.278 with Monte Carlo standard error 0.0625. Coverage nevertheless
appears to be 95.2\%, because mean width is 5.219. A nominal-looking
coverage fraction alone does not validate the design or exclude bias.

\subsection{Use prediction to economise on information}
Repeated weighting expends information on action randomness as well as
reward uncertainty. A model of future reward can remove part of this
variation. The sequential doubly robust construction of
\citet{jiang2016}, equation~(10) and Theorem~1, makes this distinction
precise. Here the model is fitted on a separate pilot and action
probabilities remain known. Unbiasedness does not require the fitted
outcome model to be correct.

For fixed pilot regressions $\widehat Q_t(H_t,A_t)$, put
$\widehat V_t(H_t)=\widehat Q_t\{H_t,\pi(H_t)\}$ and
$\widehat V_3=0$. The full-history score is
\begin{equation}
D(\pi)=\widehat V_0(H_0)+
 \sum_{t=0}^{2}W_t(\pi)
 \{R_t+\widehat V_{t+1}(H_{t+1})-\widehat Q_t(H_t,A_t)\}.
\label{closure30:sequential-score}
\end{equation}
Condition on the pilot. With $W_{-1}=1$, random assignment gives
\[
E\{W_t\widehat Q_t(H_t,A_t)\}
  =E\{W_{t-1}\widehat V_t(H_t)\}.
\]
Model terms cancel in expectation, leaving
$E(D\mid\text{pilot})=V(\pi)$. This telescoping argument covers the present
reward and transition dependence without importing a variance formula under
unstated assumptions. Setting every $\widehat Q_t$ to zero returns the
reward-weighted score.

Apply recovery to the residual contributions. Initial $\widehat V_0$ is
observed for everyone and belongs in $K_D$; split the residual sum into
routinely observed $K_D$ and missing $M_D$. Then $Z_D=K_D+B M_D/q$ has
conditional mean $D$ and
\[
\operatorname{Var}(Z_D\mid\text{pilot})
=\operatorname{Var}(D\mid\text{pilot})
 +(q^{-1}-1)E(M_D^2\mid\text{pilot}).
\]
Good predictions can reduce both terms. A wrong but fixed square-integrable
model does not invalidate the mean identity when the probabilities are
correct. Reusing evaluation outcomes for fitting would require a different
argument, such as a justified sample-splitting procedure.

\subsection{The precision improvement and its price}
A second experiment uses 5,000 independent pilot customers per
replication and random 25\% recovery of incomplete pilot histories.
Working backward, separate regressions for each policy use the fixed basis
$(1,L,X,A,L^2,X^2,LX,AL,AX)$. Routine transitions receive weight one and
recovered transitions weight four. Another 2,000 customers evaluate the
scores. There are 500 replications per observation mechanism, with fresh
seeds. The model basis is fixed before final evaluation.

\begin{table}[htbp]
\centering\small
\begin{tabular}{lrrrr}
\toprule
Adverse dropout & Bias & Coverage & Width & Positive lower\\
\midrule
Assumed $g_t$, assisted & 0.925 & 0.590 & 2.194 & 84.4\%\\
5\% recovery, assisted & 0.082 & 0.958 & 4.769 & 15.8\%\\
25\%, reward weighting & -0.079 & 0.962 & 6.094 & 4.4\%\\
25\%, constant model & -0.025 & 0.968 & 3.229 & 14.4\%\\
25\%, quadratic model & -0.003 & 0.964 & 2.773 & 24.0\%\\
50\%, quadratic model & 0.016 & 0.950 & 2.181 & 33.2\%\\
Full, quadratic model & 0.001 & 0.952 & 1.827 & 41.4\%\\
\bottomrule
\end{tabular}
\caption{Fixed-policy improvement with an independent pilot. Rows share
evaluation samples and recovery draws. Assisted and constant-model rows
also use the 5,000-customer pilot. Positive lower denotes a positive
interval endpoint, not adoption. Reward weighting needs no pilot.}
\end{table}

At 25\% evaluation recovery, the quadratic model reduces width from 6.094
to 2.773, or 54.5\%. The paired difference is $-3.320$, with Monte Carlo
standard error 0.038. Bias is $-0.0026$ with Monte Carlo standard error
0.0323. Coverage of 96.4\% has Monte Carlo standard error 0.83 percentage
points. Under conditionally independent routine dropout, coverage is 95.2\%
and width 2.607. The crude constant model also remains centred within
simulation error, while its wider intervals show the cost of weaker
prediction.

The pilot is additional information. Under adverse dropout it needs about
775 completed recoveries on average, besides roughly 309 evaluation
recoveries at 25\%. The unassisted comparator needs only evaluation
recoveries. This is not an equal-total-budget comparison. Reusing a pilot
across policies might spread its cost, but that benefit must be measured.

\begin{figure}[htbp]
\centering
\includegraphics[width=\textwidth]{twelve-hour-closure/2026-10-29/outcome-figures/precision.pdf}
\caption{More recovery narrows intervals, but positive evidence for a modest
gain remains scarce. Outcome-dependent dropout; 500 replications per study.
The assisted curve needs about 775 extra pilot recoveries, excluded from the
horizontal evaluation count. Curves use separate studies; the
paired 25\% comparison is in the text.}
\end{figure}

The false routine observation model with the assisted score produces a more
dangerous result: bias 0.925 and coverage 59.0\%, while 84.4\% of intervals
have a positive lower endpoint. Prediction sharpens a selected view of the
customers without correcting the observation mechanism.

\begin{figure}[htbp]
\centering
\includegraphics[width=\textwidth]{twelve-hour-closure/2026-10-29/outcome-figures/bias-coverage.pdf}
\caption{The observation design controls bias; regression controls part of
the variance. Bars are ordinary 95\% Monte Carlo intervals across 500
replications, not simultaneous qualification bounds. Reference uncertainty
is separate. Full history is a diagnostic benchmark, not bank data.}
\end{figure}

\subsection{Compare the same collection budget}
The preceding precision gain uses extra pilot information. A fair resource
comparison gives every method the same 7,000 initial records and the exact
same randomly recovered histories, at a 25\% recovery rate. In each of 500
new replications, one method weights all records directly. Three others
reserve 5,000, 2,000 or 1,000 customers for fitting and use the remaining
customers for evaluation. The partitions are fixed before outcomes are
examined. All use the same quadratic regression basis and fixed policies.

A fifth method divides the cohort into two equal halves. It fits a model
on each half, evaluates that model on the other half, and averages all
held-out scores. Each customer contributes to evaluation, but its own
outcome is never used to fit the model applied to it. This cross-fitting
has a scientific condition, not just a computational advantage.

\begin{table}[htbp]
\centering\small
\begin{tabular}{lrrrr}
\toprule
Adverse dropout & Bias & Coverage & Width & Positive lower\\
\midrule
Weight all 7,000 & 0.049 & 0.960 & 3.272 & 14.0\%\\
Pilot 5,000; evaluate 2,000 & 0.034 & 0.968 & 2.713 & 25.4\%\\
Pilot 2,000; evaluate 5,000 & -0.011 & 0.946 & 1.781 & 41.4\%\\
Pilot 1,000; evaluate 6,000 & 0.003 & 0.952 & 1.654 & 47.6\%\\
Two-fold cross-fitting & -0.010 & 0.944 & 1.512 & 53.8\%\\
\bottomrule
\end{tabular}
\caption{Equal realised collection resources: every method receives the same
7,000 initial records and completed recoveries, averaging 1,086 histories
under adverse dropout. There are 500 new replications. These are selected
fixed allocations, not a search proving a globally optimal split.}
\end{table}

Reducing the pilot to 1,000 customers improves precision because more
customers then evaluate a model that is already useful. Cross-fitting has
the narrowest average interval among these declared alternatives. Relative
to weighting all records, its paired reduction in width is 1.761, with
Monte Carlo standard error 0.013. Estimated bias is $-0.0103$ with Monte
Carlo standard error 0.0181. Under conditionally independent routine
dropout, its width is 1.441 and coverage is 96.6\%; every method then uses
about 895 recovered histories on average.

The adverse-case coverage of 94.4\% has Monte Carlo standard error about
1.03 percentage points. Estimated variance averages 0.154, compared with
the empirical variance of estimates, 0.163. These finite-sample findings
remain visible rather than being rounded into a calibration guarantee.
Equal collection resources also do not mean identical operating costs:
fitting, administration and recovery quality still require measurement.

\subsection{Why cross-fitting has a conditional justification here}
Let $Z_Q$ be the recovered score for any fixed outcome regressions $Q$.
Known action and recovery probabilities give $E(Z_Q)=V(\pi)$ for every
such square-integrable $Q$. Suppose each fold's fitted model approaches a
common deterministic limit $Q^*$ in the sense
\[
\|Z_{\widehat Q_{-k}}-Z_{Q^*}\|_{L_2(P)}=o_p(1).
\]
There are a fixed number of folds, each containing a nonvanishing fraction
of the independent customers. Conditional on the other folds, the mean
difference on fold $k$ has mean zero and variance
$\|Z_{\widehat Q_{-k}}-Z_{Q^*}\|_{L_2(P)}^2/n_k$.
Conditional Chebyshev's inequality makes that difference
$o_p(n^{-1/2})$. The finitely many fold remainders need not be independent
for their weighted sum to be negligible. Thus
\[
\widehat V_{\rm cf}(\pi)
=\frac{1}{n}\sum_{i=1}^{n} Z_{Q^*}(O_i)+o_p(n^{-1/2}).
\]
If the limiting score has finite, positive variance, the ordinary independent
sample central limit theorem applies. The same conditional argument makes
the empirical squared score differences vanish; Cauchy--Schwarz then shows
that pooled-score variance consistently estimates the limiting variance.
This justifies the normal approximation asymptotically under the stated
conditions, without claiming exact coverage in a finite sample.

The declared finite-dimensional regressions have a route to these conditions.
Action and recovery probabilities are bounded away from zero. The fixed
horizon, Gaussian states and polynomial rewards have the needed moments.
Positive state densities and both possible actions give nonsingular
population regression matrices. Weighted least squares converges backward
through the three decisions to deterministic projection coefficients, even
when the quadratic class is misspecified. Coefficient convergence and the
finite feature moments imply the score convergence above. A sample rank
check alone would not establish this population argument.

The exact mean identity is what makes this route relatively simple. If
treatment or recovery probabilities must be estimated, it no longer removes
all model error automatically. Products of estimation errors, dependent
customers, policy learning and selective completion require additional
arguments. Cross-fitting does not itself solve those problems.

\subsection{A collection budget is a design choice}
Let $r$ be the incomplete-record probability, $c_0$ the cost of an initial
record and $c_A$ the additional cost of a recovered history. Expected
evaluation cost is $n(c_0+rqc_A)$, plus pilot cost. For a fixed score let
$V_F=\operatorname{Var}(F_\Delta)$ and $B_M=E(M_\Delta^2)$. Estimator
variance is $(V_F-B_M+B_M/q)/n$.

With fixed $n$, more recovery always reduces the missing-information term.
With a fixed monetary budget, it can mean fewer initial customers. If
$A=V_F-B_M>0$, differentiating
$(A+B_M/q)(c_0+rqc_A)$ gives the unconstrained optimum
\[
q^*=\sqrt{\frac{B_Mc_0}{A r c_A}},
\]
restricted to $[q_{\min},1]$. If $A\leq0$, that expression decreases with
$q$ and favours $q=1$ within this cost model. This assumes constant costs
and a fixed trained score, not an optimal joint pilot/evaluation design.
No bank prices are supplied, so no numerical recovery fraction is recommended.

Another route should be separated from the tested regression. For information
$O$ available before recovery and a separately trained predictor $m(O)$ of
the missing contribution,
\[
Z_m=K+m(O)+\frac{B}{q(O)}\{M-m(O)\}
\]
is unbiased under the known conditional recovery design even when $m$ is
wrong. Set $m=0$ and $q=1$ for complete records, which need no recovery; the design
and allocation calculations concern incomplete records. Its extra variance is
$E[\{q(O)^{-1}-1\}\{M-m(O)\}^2]$.
The best squared-error predictor is $E(M\mid O)$. With known residual
second moment $a(O)=E[\{M-m(O)\}^2\mid O]$ and cost $c(O)>0$, minimising
this term at a fixed expected recovery budget yields
$q(O)\propto\sqrt{a(O)/c(O)}$, subject to positivity and the ceiling one.
This is a derivation, not a tested allocation algorithm. A separate pilot
must estimate the quantities; fresh data must test the design. Unseen
future outcomes cannot determine which customer is recovered.

\subsection{What this changes in the proposal}
A prototype can now be proposed around a specific design: accept declared
recovery probabilities, retain the initial customer denominator, estimate
paired policy improvements, and display the extra information burden and
remaining uncertainty. It should distinguish dropout weights from recovery
weights and show when an outcome model improves only precision. Where
recoverable histories do not exist, sensitivity or partial-identification
analysis is the appropriate next step.

These results concern independent generated customers, known action
probabilities and fixed policies. They do not settle calendar dependence,
unmeasured treatment confounding, erroneous or partially recovered histories,
estimated treatment probabilities, or uncertainty after policy learning.
Normal intervals are assessed rather than declared exact. The simulation
keeps every unaudited missing value unavailable to fitting and estimation.
The generated evidence does not establish that real histories can be
recovered accurately, or that the collection expense is worthwhile.

The remaining question is whether a feasible collection design can deliver
accurate histories in sufficient numbers to separate the economic
alternatives at an acceptable total cost.

\section{When a random invitation still leaves selective recovery}
\label{rel73:partial-recovery}

The full-history design answers the preceding question only when a selected
history can actually be recovered. A bank may instead invite a random sample
of customers whose records are incomplete and receive usable evidence from
only some of them. Customers with a resolved dispute may reply more readily
than those still in difficulty; a collections call can itself alter a later
repayment. Randomizing the invitation does not make the completed responses
representative, nor does it make a later outcome equal to the historical
outcome the bank intended to recover. Both distinctions change the conclusion
that the bank can draw from the follow-up expense.

Consider a simpler decision with a genuinely bounded endpoint. A servicing
offer is randomized before a fixed repayment date. The endpoint is whether
the specified payment was made by that date, and the target is the average
effect of the offer in the experimental customer population. The bank later
tries to recover the historical payment status from incomplete records. This
example concerns a binary historical outcome and a single treatment. It does
not impose artificial finite bounds on the earlier Gaussian servicing rewards
or claim to recover a longitudinal policy value from a terminal indicator.

\subsection{The information supplied by an invitation}

Let $A\in\{0,1\}$ denote randomized treatment and $Y(a)\in\{0,1\}$ the
historical outcome under assignment $a$. Consistency and no interference give
$Y=Y(A)$; random assignment identifies the population contrast
$\tau=E\{Y(1)-Y(0)\}$ from the two full arm means if those outcomes were
observed. Let $R$ indicate that routine records contain $Y$. For customers
with $R=0$, let $B$ indicate an invitation and let $S$ indicate whether the
specified follow-up would recover an outcome if that customer were invited.
An uninvited customer's $S$ is not observed. The design requires
\[
 B\mathbin{\perp}(Y,S)\mid A=a,R=0,\qquad
 \Pr(B=1\mid A=a,R=0)=q_a>0.
\]
The invitation probabilities are known. The outcome among completed
follow-ups must be the same historical $Y$, not a new measurement of a
changed endpoint. Neither routine response nor follow-up completion is
assumed independent of $Y$.

For each arm define the routine success mass, recovered success mass and
remaining mass by
\[
 v_a=E(RY\mid A=a),\qquad
 z_a=E\{(1-R)SY\mid A=a\},\qquad
 m_a=E\{(1-R)(1-S)\mid A=a\}.
\]
Random invitation identifies $z_a$ and $m_a$ through the observed conditional
distribution of follow-up response and outcomes among invitees. It does not
identify the outcome mean in the remaining mass. Binary support consequently
gives arm bounds $\ell_a=v_a+z_a$ and $u_a=\ell_a+m_a$, and effect bounds
\begin{equation}
 I=[\tau_L,\tau_U]=[\ell_1-u_0,\ u_1-\ell_0],\qquad
 \tau_U-\tau_L=m_1+m_0.
 \label{rel73:identified-set}
\end{equation}
This is the binary specialization of the double-sampling bounds developed
by \citet{rel73coppock2017}. Their outcome-stability assumption permits
arbitrary differences between initial respondents and nonrespondents. It
restricts the meaning of the recovered outcome, not the selectivity of the
initial response.

The bounds are attainable within the declared model. Assign outcome zero
to every unresolved treated customer and one to every unresolved control
customer to obtain the lower endpoint; reverse those assignments for the
upper endpoint. These assignments leave the routine observations,
invitation law, completion probabilities and observed follow-up outcomes
unchanged. The joint coupling of the two potential outcomes is otherwise
unrestricted, so the arm completions can coexist in a population. This
establishes sharpness relative to these observed data and assumptions; extra
restrictions on outcomes or selection would define a different model.

Holding the follow-up procedure fixed, changing a positive invitation
probability does not change $m_a$ or the population identified set. It
changes how precisely the available follow-up law can be estimated. A more
effective recovery procedure can reduce $m_a$ and hence identification width.
At $q_a=0$, the follow-up law is not identified in that arm, so the conclusion
cannot be extrapolated to no invitations. If increasing invitations changes
staff attention, response behavior or outcomes, it also changes the procedure;
the fixed-law comparison no longer applies.

\subsection{Estimating a set and covering it are separate operations}

For independent customer observations across and within the two arms,
and invitations drawn independently across customers, known Bernoulli
invitation permits
the following observable scores, with missing products evaluated only for
records whose outcomes are observed:
\begin{equation}
 L=RY+\frac{(1-R)BSY}{q_a},\qquad
 U=RY+(1-R)\left\{1-\frac{BS(1-Y)}{q_a}\right\}.
 \label{rel73:endpoint-scores}
\end{equation}
Conditional expectation over $B$ gives $E(L\mid A=a)=\ell_a$ and
$E(U\mid A=a)=u_a$. Thus
$\widehat\tau_L=\overline L_1-\overline U_0$ and
$\widehat\tau_U=\overline U_1-\overline L_0$ are unbiased endpoint
estimators. These are authored inverse-invitation estimators under Bernoulli
sampling. They are not a native execution of the published package's
sample-ratio estimators or its fixed-follow-up-size variance calculation.

Inverse weighting can give an individual score outside the outcome support.
Here $L\in[0,1/q_a]$ and $U\in[1-1/q_a,1]$, so each range has length
$1/q_a$. In a finite sample the two estimated endpoints can cross. Sorting
them or concealing such a sample would change the declared procedure. The
raw estimates therefore remain in the results, along with any empty
confidence set.

A conservative finite-sample guarantee follows directly from the score
ranges. Let $n_a$ be the number of independent customers in arm $a$, and put
\begin{equation}
 D=\frac{1}{n_1q_1^2}+\frac{1}{n_0q_0^2},\qquad
 t_\alpha=\sqrt{\frac{D}{2}\log\frac{2}{\alpha}},\qquad
 C_\alpha=[\widehat\tau_L-t_\alpha,
            \widehat\tau_U+t_\alpha]\cap[-1,1].
 \label{rel73:confidence-set}
\end{equation}
Under the stated independence, design and stability assumptions,
$\Pr(I\subseteq C_\alpha)\geq1-\alpha$. The guarantee covers the entire
population identified interval, and consequently the true effect within it.
It does not say that the reported set is narrow enough to support an action.

For completeness, the argument needs no estimated variance. If a random
variable has range length $c$, the second derivative of its centered
log-moment-generating function is its variance under exponential tilting.
That variance is at most $c^2/4$: the squared distance from the range
midpoint is bounded by $c^2/4$, and the variance is no larger. Integrating
twice from zero gives a log-moment bound of $s^2c^2/8$. Independence then
bounds either centered endpoint estimate's log moment by $s^2D/8$.
Exponential Markov inequality, minimized at $s=4t/D$, bounds either adverse
tail by $\exp(-2t^2/D)$. Applying this to overestimation of the lower endpoint
and underestimation of the upper endpoint, and summing the two failure
probabilities, proves~\eqref{rel73:confidence-set}. Dependence between the
two endpoint estimates is allowed; independence of the customers used in
the product argument is not optional. Intersection with $[-1,1]$ preserves
coverage because the population identified interval is inside that support.

The comparison also reports normal endpoint bands using sample variances of
the scores and a 1.96 multiplier on each outward standard error. Their
behavior is empirical here; the preceding finite-sample guarantee does not
transfer to them. The original article's Imbens--Manski construction instead
targets asymptotic coverage of the true effect, which is a different
criterion from simultaneous containment of both endpoints. A software label
such as \emph{confidence interval around the region} does not erase this
distinction. \citep{rel73coppock2017,rel73attrition}

\subsection{Selective completion can reverse a point conclusion}

The retained generated experiment has historical arm means .45 and .55,
so the true effect is .10. Routine observation probabilities, indexed by
outcome zero and one, are $(.8,.4)$ in control and $(.4,.8)$ in treatment.
Both arms therefore have routine response probability .62, but their observed
success means differ sharply. Equal response rates do not imply equal
selection. The routine-only complete-case contrast is approximately .4194,
while routine-only worst-case bounds are $[-.12,.64]$.

Three stable recovery procedures vary successful completion among invited
nonresponders. The standard procedure has probabilities $(.1,.9)$ in control
and $(.9,.2)$ in treatment; the intensive procedure has $(.8,.98)$ and
$(.98,.8)$ respectively; complete recovery has probability one in both
arms. These probabilities define the generated truth. The estimator is not
given the outcome-dependent selection law or the true arm means.

\begin{table}[htbp]
\centering\small
\begin{tabular}{lrrr}
\toprule
Recovery procedure & Population effect set & Width & Naive limit\\
\midrule
Routine records only & $[-.1200,.6400]$ & .7600 & .4194\\
Standard & $[-.0870,.1540]$ & .2410 & $-.0720$\\
Intensive & $[.0560,.1108]$ & .0548 & .0697\\
Complete & $[.1000,.1000]$ & .0000 & .1000\\
\bottomrule
\end{tabular}
\caption{Exact generated population quantities. The routine-only naive
contrast uses complete cases. The other naive limits replace every routine
nonrespondent's outcome by the mean among completed follow-ups in that arm.
All stable sets contain the true effect .10.}
\label{rel73:population}
\end{table}

Under standard recovery, successful control follow-ups overwhelmingly reveal
historical successes, whereas successful treated follow-ups disproportionately
reveal failures. Treating those replies as representative produces the
negative limiting contrast in Table~\ref{rel73:population}. A random
invitation has not repaired selective completion. The valid set retains
both signs; the intensive procedure supplies enough extra outcome information
to exclude zero at the population level.

Each procedure is evaluated in 500 independent datasets with 2,000 customers
per arm. A fourth, deliberately unstable procedure is discussed below.
Within a dataset, shared invitation uniforms compare probabilities
.10, .25, .50 and 1.00. There are 2,000 independent generated datasets and
8,000 procedure--invitation-rate cells, not 8,000 independent datasets.
All cells remain in the denominator. Reaggregation checks every retained
endpoint, interval indicator, resource total and summary; direct observed-cell
sums and squares reconstruct endpoint estimates and standard errors from
the first, middle and last retained raw datasets in each procedure.

\begin{table}[htbp]
\centering\small
\begin{tabular}{llrrrr}
\toprule
Procedure & Invitation & \multicolumn{2}{c}{Whole-set coverage}
 & \multicolumn{2}{c}{Mean set width}\\
 & probability & Bounded & Normal & Bounded & Normal\\
\midrule
Standard & .10 & 100.0\% & 97.6\% & 1.100 & .376\\
Standard & 1.00 & 99.2\% & 96.0\% & .327 & .302\\
Intensive & .10 & 100.0\% & 96.6\% & .915 & .220\\
Intensive & 1.00 & 99.4\% & 94.8\% & .141 & .117\\
Complete & .10 & 100.0\% & 98.0\% & .860 & .172\\
Complete & 1.00 & 99.2\% & 95.4\% & .086 & .062\\
\bottomrule
\end{tabular}
\caption{Selected endpoints of the retained invitation comparison, with
500 replications per row and nominal 95\% confidence. Bounded refers to
equation~\eqref{rel73:confidence-set}; Normal is the separately assessed
normal band. Coverage requires containment of the entire population set.
An empty confidence set has width zero and fails coverage.}
\label{rel73:coverage}
\end{table}

The bounded construction buys a finite-sample statement with substantial
conservatism. Under intensive recovery at invitation probability .10 it
declares no positive effects, even though the population lower endpoint is
.056. At probability one it returns a nonempty strictly positive set in
401 of 500 datasets. The corresponding normal-band counts are 193 and
469. Mean invitations rise from about 152 to 1,521, and mean completed
recoveries from about 141 to 1,411. These counts describe information burden,
not bank collection prices or the superiority of either interval rule.

At 94.8\% coverage the Monte Carlo standard error is approximately one
percentage point. The difference from nominal 95\% therefore does not by
itself establish miscalibration; the theorem and its absence for the normal
band remain separate from these empirical proportions. Across all four
invitation rates, raw endpoint crossings occur in 108 intensive-recovery
cells and 712 complete-recovery cells. One complete-recovery normal band is
empty. None is removed from coverage counts, and no empty set is treated
as a supported positive decision. The original signed endpoint separations
remain available alongside the nonnegative set widths used here.

\subsection{A recovered record can concern the wrong outcome}

The fourth procedure completes every invited follow-up but changes the
recorded endpoint for half of the originally unsuccessful, routinely missing
treated customers. It leaves control outcomes unchanged. This deliberately
violates historical-outcome stability. The apparent identified point becomes
.235 although the historical treatment effect remains .10. Inviting more
customers now estimates the wrong contrast more precisely.

At invitation probabilities .10, .25, .50 and 1.00, the bounded sets cover
the historical .10 in 100.0\%, 90.4\%, 0.2\% and 0.0\% of the 500
datasets. The normal bands cover it in 7.0\%, 0.2\%, 0.0\% and 0.0\%.
The wide first band's accidental coverage is not evidence that the stability
assumption is dispensable. At full invitation the bounded construction
covers the different recorded-outcome point .235 in all 500 datasets,
while the normal band covers it in 472. Neither statement validates
inference for the historical target.

A sensitivity analysis must restrict the error that actually matters.
Let $k_a=E\{(1-R)S\mid A=a\}$ and, for $k_a>0$, let
$\widetilde\mu_a$ be the mean recorded outcome among recoverable customers.
A zero recovered mass contributes zero without defining that mean.
Suppose external information
justifies an absolute bound $\gamma_a$ on the difference between that mean
and their historical-outcome mean. The recovered contribution then lies in
\[
 k_a\max(0,\widetilde\mu_a-\gamma_a)
 \ \leq\ E\{(1-R)SY\mid A=a\}\ \leq\
 k_a\min(1,\widetilde\mu_a+\gamma_a).
\]
Adding the routine contribution and unresolved mass gives a sensitivity set
by the same arm subtraction as~\eqref{rel73:identified-set}. This is a
population sensitivity calculation. Estimating its endpoints would require
an additional sampling argument; the original confidence formula cannot be
attached without checking the altered scores and assumptions.

In the unstable construction, $k_1=.38$, the recovered treated success mass
is .245, and its historical value is .11. If control recovery is stable,
the smallest treated mean-drift bound that admits the truth is
$\gamma_1=.135/.38=27/76$, approximately .3553. At $\gamma_1=.30$ the
effect sensitivity set is $[.121,.349]$ and still excludes .10; at .40 it
is $[.083,.370]$. These are stipulated limits on outcome drift, not estimates
learned from the incomplete data.

The distinction matters when choosing software. The inspected
\texttt{attrition} 1.0.0 source at its 21 September 2026 revision implements
the paper's double-sampling bounds, analytic variance and normal critical-value
calculation. Its package description calls the sensitivity analysis a check
on outcome stability, but the function documentation and
\texttt{gen\_mean} implement a different restriction: a fraction $1-\delta$
of follow-up nonresponders has the responder mean, while the remaining
fraction $\delta$ receives worst-case outcomes. That is the paper's
nonresponse-mixture sensitivity, not a bound on historical-outcome drift.
When recovery is complete, there are no follow-up nonresponders and changing
$\delta$ cannot alter the erroneous .235 contrast. Static inspection
establishes this distinction; it is not a native package trial.
\citep{rel73attrition,rel73coppock2017}

\subsection{The information purchase and its alternatives}

The available implementation is a serious comparator for bounded randomized
experiments with the declared follow-up design. Its sample-ratio estimator,
asymptotic variance and poststratification should be compared under matched
sampling rules and targets. The proposed inverse-invitation calculation makes
known design probabilities and a conservative finite-sample guarantee explicit,
but its crossed endpoints and wide intervals are costs of that choice. Neither
route validates the other merely because their population limits agree.
The package declares GPL-3; a proposed integration must retain its version,
applicable terms and dependency record. No enabled integration or measured
runtime advantage is claimed. \citep{rel73attrition}

Other options change the information or assumptions. A reliable administrative
linkage may recover the fixed historical outcome more credibly than a later
customer report. An intensive protocol can reduce the unresolved mass, while
more invitations to the same protocol mainly improve precision. A model that
assumes completion is ignorable conditional on recorded covariates can support
a point estimate, but it must justify that assumption and diagnose overlap.
The package's trimming option instead targets customers whose outcomes would
be observed under either treatment, with selection assumptions explicitly
chosen. That always-observed target can be useful, but it is not the full
experimental-population effect in this section. Equal observed response rates
alone do not establish individual response monotonicity or justify changing
the target without saying so. \citep{rel73attrition}

To connect identification with a financial decision, stipulate a benefit of
USD~100 for each additional historical repayment and an offer cost of USD~5
per customer. These invented values define net benefit $100\tau-5$, not
observed bank returns. The standard-recovery population set permits net
values from USD~$-13.70$ to USD~10.40. Intensive recovery narrows this to
USD~0.60 through USD~6.08; complete stable recovery gives USD~5.00. Those
are population implications. A finite-sample spending recommendation needs
its own nonempty net-benefit confidence set above zero, and follow-up costs
must also be counted. A positive effect alone is insufficient to clear the
USD~5 threshold.

At full invitation under intensive recovery, only 5 of 500 bounded sets
and 22 normal bands have a nonempty net-benefit range wholly above zero.
The earlier counts of 401 and 469 concerned a positive effect, before
deducting the offer cost. Under complete stable recovery, the corresponding
net-positive counts are 335 and 432. These are decision frequencies in the
same retained datasets. They show how a small population margin can remain
hard to distinguish from zero even after every missing customer is invited. The normal-band frequencies retain their empirical qualification;
more affirmative decisions alone do not establish a better procedure.

The resulting design choice is concrete. The bank should establish whether
follow-up recovers the fixed historical endpoint, record every invitation and
completion, and compare the cost of reducing the unresolved mass with the
cost of estimating a fixed set more precisely. A workflow that cannot defend
stability should expose an appropriately justified sensitivity range or
withhold the historical-effect recommendation. These generated results close
the missing selective-completion argument for a bounded single-treatment
example. They do not establish feasible bank recovery, a general solution
for unbounded longitudinal rewards, or independent operating acceptance.


